\section{实变函数概要}\label{sec:Real_Analysis}

读者可能已经注意到，并不是所有的函数都是可积的，例如$\mathrm{e}^x$即使视为功率信号，其能量和功率也不具有意义，因此它既不属于能量信号，也不属于功率信号；另一种情况下，函数本身的性质比较差，以至于在黎曼积分的框架下说不清它的积分（或能量）是什么：
\begin{example}[狄利克雷函数]
    \[\chi_{\mathbb{Q}}(x)=
    \begin{cases}
        1 & x \in \mathbb{Q}    \\
        0 & x \in \mathbb{Q} ^c
    \end{cases}\]
（其中$\mathbb{Q}$ 表示有理数集，上标c表示取补集）
\end{example}
它在黎曼积分的意义下是不可积的，因为它在所有点不连续，但是这样的函数在\textbf{勒贝格积分}的意义下可积。在勒贝格积分的框架下，有理数集这样的可数集是零测度的（类似于区间长度为0），$\chi_{\mathbb{Q}}$在零测集上取1而在其余位置取0，所以其积分值为0。

对于非反常积分，\textbf{勒贝格积分是黎曼积分的推广}，黎曼可积的函数一定是勒贝格可积的，并且在相同的区间上积分值相等。

本节就来介绍一些基础的实变函数理论，以便后续学习中使用。由于篇幅原因，我们不得不略去大多数证明（证明可以参考\cite{real_analysis}），不过读者不妨先接受其中的一些重要结论，尤其是有黑体字阐释的定理。

\subsection{*勒贝格测度与可测函数}\label{sub:lebesgue}

首先，我们给出\textbf{外测度}和\textbf{可测集}的概念。
\begin{definition}[勒贝格外测度]
    设集合$A\in\mathbb{R}$，其勒贝格外测度定义为
    \[m^*(A)=\inf\left\{\sum_{k=1}^{\infty}l(I_k)\left|A\subset\bigcup_{k=1}^{\infty}I_k\right.\right\}\]
    其中$\{I_k\}_{k=1}^{\infty}$为覆盖$A$的可数开区间族，$l(I_k)$表示区间$I_k$的长度。
\end{definition}
这个定义看起来不是很直观，但它很好地\textbf{描述了“集合的长度”这一概念}，例如，对于任何一个区间$I$来说，总有$l(I)=m^*(I)$。此外，外测度的以下基本性质符合我们的直观：
\begin{proposition}[外测度的性质]
    \quad
    \begin{itemize}
        \item 单调性：设$A\subset B$，则$m^*(A)\leq m^*(B)$
        \item 平移不变性：实数集$A$的平移记为为$A+y=\{x\in \mathbb{R}|x-y\in A\}$，满足$m^*(A+y)=m^*(A)$
        \item 次可数可加性：设$\{E_k\}_{k=1}^{\infty}$是可数集族，则
        $$m^*\left(\bigcup_{k=1}^{\infty}E_k\right)\leq\sum_{k=1}^{\infty}m^*(E_k)$$
    \end{itemize}
\end{proposition}

次可数可加性中不等号的提出是为了处理$\{E_k\}_{k=1}^{\infty}$之间有重叠的情况，但是一般而言，即使它是不交集族，也不满足可数可加性，即
$$m^*\left(\bigsqcup_{k=1}^{\infty}E_k\right)\neq\sum_{k=1}^{\infty}m^*(E_k)$$
其中$\sqcup E_k$表示集族$\{E_k\}$是彼此不交的。保证这样的可数可加性，需要集合列中的每个集合\textbf{可测}。

\begin{definition}[可测集]
    集合$E$是\textbf{可测集}，如果
    $$\forall A,m^*(A)=m^*(A\cap E)+m^*(A\cap E^c)$$
    这个条件称为称\textbf{Carathéodory条件}或\textbf{卡氏条件}。定义可测集的\textbf{勒贝格测度}为它的勒贝格外测度，即
    $m(E)=m^*(E)$。
\end{definition}
在这个约束下，集合的不交族的勒贝格测度是可数可加的，即：
\[m\left(\bigsqcup_{k=1}^{\infty}E_k\right)=\sum_{k=1}^{\infty}m(E_k)\]
此外，区间都是可测的，它们的测度就等于区间的长度。我们所遇到的几乎所有集合都是可测的，但是，根据Vitali定理\cite{real_analysis}，任意具有正测度的集合，都有不可测子集，这个子集的构造可以借助选择公理给出。

\begin{proposition}
    记$\mathbb{R}$上的可测集族为$\mathcal{M} $，则\\
    \ding{172}空集$\varnothing $，全集$\mathbb{R}$均属于$\mathcal{M} $；\\
    \ding{173}若$E\in \mathcal{M} $，则$E^c\in \mathcal{M} $；\\
    \ding{174}可测集的可数族之并仍然是可测集，即
    $$\{E_k\}_{k=1}^{\infty}\subset \mathcal{M}\implies\bigcup_{k=1}^{\infty}E_k\in \mathcal{M} $$
    因此，$\mathcal{M} $构成$\mathbb{R}$上的$\sigma$-代数（我们不去定义这个概念，仅仅指出\textbf{可测集在常规的集合运算下是封闭的}）。通过以上三条性质可以推知：可测集的可数族的交仍然是可测集；两个可测集的差集仍然是可测集。
\end{proposition}

\begin{proposition}
    外测度为0的集合都是可测的，称之为\textbf{零测集}。
\end{proposition}
\begin{proof}
    设$m^*(E)=0$，则对于任意集合$A$，$A\cap E,A\cup E^c\subset E$，根据外测度的单调性，有$m^*(A\cap E),m^*(A\cap E^c)\leq m^*(E)=0$，因此$m^*(A\cap E)=m^*(A\cap E^c)=0$。
    因此，
    \[m^*(A)=m^*(A\cap E)+m^*(A\cap E^c)=0\]
\end{proof}

\begin{example}
    单点集是零测集。由测度的可数可加性，可数集是零测集，特别地，$\mathbb{Q}$是可数的，因而是零测集。但是，零测集未必是可数的，对于这样的例子，可以参考\cite{real_analysis}中的康托集，它可以视为不可数单点集之并，但其测度为0。
\end{example}

\begin{definition}[几乎处处]
    对于可测集$E$，称一个性质在$E$上\textbf{几乎处处}(almost everywhere, a.e.)成立，如果存在一个零测子集$E_0$，使得该性质对所有$x\in E\setminus E_0$成立。
\end{definition}

\begin{example}[几乎处处相等]
    函数$f,g:E\to\mathbb{R}$在$E$上\textbf{几乎处处相等}，如果存在一个零测子集$E_0$，使得对于所有$x\in E\setminus E_0$，有$f(x)=g(x)$，记为$f=g,\text{a.e.}$或$f\sim g$。此时，我们\textbf{不区分$f$和$g$}，因为这样的两个函数性质类似，并且我们将看到$f,g$可积时积分值相等。
\end{example}

下面我们引入\textbf{可测函数}的概念。

\begin{definition}[可测函数]
    定义在可测集上的函数$f:E\to\mathbb{R}$是可测的，如果
    $$\forall c\in\mathbb{R},E_c=\left\{x\in E\left| f(x)<c\right.\right\}\text{可测}$$
    可测函数构成线性空间。
\end{definition}
\begin{example}
    基本初等函数在其定义域上都是可测函数。
\end{example}

这些公理化的定义为理论推带来了方便，但对初学者而言往往难以理解。下面给出可测集和可测函数的一些刻画，以帮助读者快速建立对它们的直观。数学家Littlewood曾说\cite{real_analysis}，只要掌握了以下三个观点，就相当于掌握了可测集与可测函数的核心内容。

首先，\textbf{可测集可以用区间的有限不交并近似}：

\begin{theorem}
    设$E$可测，$m(E)<\infty$，则$\forall \epsilon>0$，存在开有界区间族$\{I_k\}_{k=1}^{n}$，记$U=\bigsqcup_{k=1}^{n} I_k$，有
    \[m(U\setminus E)+m(E\setminus U)<\epsilon\]
\end{theorem}

其次，\textbf{可测函数接近于连续函数}：

\begin{theorem}[Lusin定理]
    设$f:E\to\mathbb{R}$可测，则$\forall \epsilon>0$，存在连续函数$g:E\to\mathbb{R}$及闭集$F\subset E$，使得$m(E\setminus F)<\epsilon$，且在$F$上有$f=g$。
\end{theorem}

在勒贝格测度的框架下，逐点收敛性可以减弱为\textbf{几乎处处收敛性}：

\begin{definition}[几乎处处收敛]
    函数列$\{f_n:E\to\mathbb{R}\}_{n=1}^{\infty}$在$E$上几乎处处收敛于函数$f:E\to\mathbb{R}$，如果存在一个零测度集$E_0$，使得
    $$\forall x\in E\setminus E_0,\lim_{n\to\infty}f_n(x)=f(x)$$
    记为$f_n\to f,\text{a.e.}$
\end{definition}

\begin{example}
    设$E=[0,1]$，函数列$\{x^n\}_{n=1}^{\infty}$在$E$上逐点收敛于函数
\[f(x)=
\begin{cases}
0 & x\in [0,1)    \\
1 & x=1
\end{cases}\]
因此$\{x^n\}$几乎处处收敛于0.
\end{example}

对于可测函数，若不区分几乎处处相等的函数，则几乎处处收敛就相当于逐点收敛。在这样的约定下，函数列\textbf{在零测集上不好的性质几乎不会影响到全局性质}。

Littlewood的第三个观点是，\textbf{几乎处处收敛的可测函数列近似于一致收敛}：

\begin{theorem}[Egoroff定理]
    设$m(E)<\infty$，可测函数列$\{f_n:E\to\mathbb{R}\}$几乎处处收敛于$f:E\to\mathbb{R}$，则$\forall \epsilon>0$，存在闭集$F\subset E$，满足$m(E\setminus F)<\epsilon$，且在$F$上$f_n$一致收敛于$f$，记为$f_n\rightrightarrows f$。
\end{theorem}

\subsection{*勒贝格积分}\label{sub:integral}

我们首先回顾黎曼积分的定义。

\begin{definition}[黎曼积分]
    设$f:[a,b]\to\mathbb{R}$有界，$P=\{x_0,x_1,\dots,x_n\}$是区间$[a,b]$的一个划分，$a=x_0<x_1<\cdots <x_n=b$，对于每个子区间$[x_{k-1},x_k]$，定义
    \[m_k=\inf_{x\in [x_{k-1},x_k]}f(x),M_k=\sup_{x\in [x_{k-1},x_k]}f(x)\]
    由此定义$f$对$P$的上、下\textbf{达布和}(Darboux sum)分别为
    \[U(f,P)=\sum_{k=1}^{n}M_k(x_k-x_{k-1}),L(f,P)=\sum_{k=1}^{n}m_k(x_k-x_{k-1})\]
    如果
    \[\inf_{P}U(f,P)=\sup_{P}L(f,P)\]
    则称$f$在$[a,b]$上\textbf{黎曼可积}，并将上述公共值定义为$f$在$[a,b]$上的\textbf{黎曼积分}，为了区分于勒贝格积分，记为$(R)\int_{a}^{b}f(x)\,dx$。
\end{definition}

这个定义实质上实质上是先定义了\textbf{阶梯函数}的黎曼积分，然后通过阶梯函数对被积函数进行上、下逼近，从而定义了被积函数的黎曼积分。阶梯函数的定义如下：

\begin{definition}[阶梯函数]
    如果函数$\varphi:[a,b]\to\mathbb{R}$的定义域可以被划分为有限的子区间族$\{[x_{k-1},x_k]\}_{k=1}^{n}$（严格来说需要处理区间端点，但区间端点属于左侧、右侧的集合均没有影响），使得$\varphi$在每个子区间的内部$(x_{k-1},x_k)$为常数函数，则称$\varphi$是定义在$[a,b]$上的\textbf{阶梯函数}。

    显然阶梯函数黎曼可积，并且
    $$\int_{a}^{b}\varphi(x)\,dx=\sum_{k=1}^{n}\varphi(x)(x_k-x_{k-1})$$
\end{definition}

勒贝格定理\cite{real_analysis}表明，闭区间上的有界函数$f$黎曼可积的充要条件是$f$在$[a,b]$中的不连续点构成零测集，这对函数的要求是比较高的，例如狄利克雷函数$\chi_{\mathbb{Q}}$在$[0,1]$不连续，因此它在黎曼积分的意义下不可积，而\textbf{勒贝格积分扩展了有界区间上可积函数的范围}，这是勒贝格积分的第一条重要优势。

勒贝格积分的定义与黎曼积分类似，区别在于它需要通过\textbf{简单函数}来进行逼近。

\begin{definition}[简单函数]
    若存在函数$\varphi:E\to\mathbb{R}$的定义域$E$的一个可测集划分$\{E_k\}_{k=1}^{n}$，使得$\varphi$在每个可测子集$E_k$上为常数函数，则称$\varphi$是定义在$E$上的\textbf{简单函数}。

    简单函数有\textbf{典范表示}(canonical representation)，即
    \[\varphi=\sum_{k=1}^{n}a_k\chi_{E_k},\text{其中}E_k=\varphi^{-1}(a_k)=\{x\in E|\varphi(x)=a_k\}\]
\end{definition}

显然，简单函数是可测的，并且它的勒贝格积分应该定义为
\[\int_{E}\varphi\,dm=\sum_{k=1}^{n}a_k m(E_k)\]
$dm$表示勒贝格测度。
\begin{remark}*
    对于单调递增的有界变差函数$f$，可以由$f$的增量定义新的测度，并由此引出\textbf{勒贝格-斯蒂尔杰斯积分}，此时$dm$就应替换为$df$，但我们不在这里展开讨论，感兴趣的读者可以参考\cite{simple_real_analysis}或\cite{chentianquan}，它实际上是分部积分的基础。
\end{remark}

下面，我们先通过简单函数定义有界函数的勒贝格积分，然后扩展被积函数的范围，定义非负可测函数的勒贝格积分，再由此引出任意可测函数的勒贝格积分。

\begin{definition}[有界函数的勒贝格积分]
    设$f:E\to\mathbb{R}$是有界可测函数，如果任给简单函数$\varphi\leq f$和$\psi\geq f$，$\varphi,\psi$仅在有限测度集上不为0，都有
    \[\sup_{\varphi\leq f}\int_{E}\varphi\,dm=\inf_{\psi\geq f}\int_{E}\psi\,dm\]
    则称$f$在$E$上\textbf{勒贝格可积}，并将上述公共值定义为$f$在$E$上的\textbf{勒贝格积分}，记为$\int_{E}f\,dm$。
\end{definition}

\begin{theorem}
    设$f:E\to\mathbb{R}$是有界可测函数，则$f$在$E$上勒贝格可积；反之，设$m(E)<\infty$，如果有界函数$f$在$E$上勒贝格可积，则$f$在$E$上可测。
\end{theorem}

这个定理表明，可测函数是勒贝格积分理论中的主要研究对象，在多数情况下，$f$可测就等价于能够指出$f$的勒贝格积分的值，并且$m(E)<\infty$时$f$还是可积的。

\begin{theorem}
    设$f:[a,b]\to\mathbb{R}$是黎曼可积函数，则$f$也是勒贝格可积的，并且在$[a,b]$上的黎曼积分与勒贝格积分相等，即
    \[(R)\int_{a}^{b}f(x)\,dx=\int_{a}^{b}f\,dm\]
\end{theorem}

在这个意义下，\textbf{勒贝格积分是黎曼积分的推广}。我们将在后文中明确区间上的勒贝格积分记为$\int_{a}^{b}f\,dm$的具体意义，但这样的记号是自然的。此外，勒贝格积分具有如下性质：
\begin{proposition}[有界可测函数的勒贝格积分的线性性]
    设$f,g:E\to\mathbb{R}$是有界可测函数，$\alpha,\beta\in\mathbb{R}$，如果$f,g$在$E$上勒贝格可积，则$\forall \alpha,\beta\in\mathbb{R}$，$\alpha f+\beta g$在$E$上勒贝格可积，并且
    \[\int_{E}(\alpha f+\beta g)\,dm=\alpha\int_{E}f\,dm+\beta\int_{E}g\,dm\]
\end{proposition}
\begin{corollary}[有界可测函数的勒贝格积分的单调性]
    设$f,g:E\to\mathbb{R}$是有界可测函数，如果$f,g$在$E$上勒贝格可积，并且对于所有$x\in E$，有$f(x)\leq g(x)$，则
    \[\int_{E}f\,dm\leq \int_{E}g\,dm\]
\end{corollary}
\begin{corollary}[有界可测函数的勒贝格积分的区域可加性]
    设$f:E\to\mathbb{R}$是有界可测函数，$A,B\subset E$是可测不交子集，则
    \[\int_{A\sqcup B}f\,dm=\int_{A}f\,dm+\int_{B}f\,dm\]
\end{corollary}

勒贝格积分向误解函数的推广需要先定义非负可测函数的勒贝格积分：

\begin{definition}[非负可测函数的勒贝格积分]
    设$f:E\to[0,+\infty)$是非负可测函数，则定义$f$在$E$上的\textbf{勒贝格积分}为
    \[\int_{E}f\,dm=\sup_{0\leq h\leq f}\int_{E}h\,dm\]
    其中$h$为定义在$E$上的有界可测函数，并且仅在有限测度集上不为0。显然非负可测函数的勒贝格积分是非负的。
    
    称$f$是\textbf{可积}的，如果$\int_{E}f\,dm<\infty$。注意不可积的函数，也可指出其勒贝格积分的值为$+\infty$。
\end{definition}

如果$f$还是有界函数，则以上定义与有界函数的勒贝格积分定义是一致的，因此是良定义的。

\begin{theorem}[非负可测函数的勒贝格积分的线性性]
    设$f,g:E\to[0,+\infty)$是非负可测函数，$\forall\alpha,\beta\in[0,+\infty)$，有
    \[\int_{E}(\alpha f+\beta g)\,dm=\alpha\int_{E}f\,dm+\beta\int_{E}g\,dm\]
\end{theorem}

它的证明相对简单，且有助于理解“简单函数的下逼近”这一概念，因此我们在这里给出证明。
\begin{proof}
    不妨设$\alpha=\beta=1$，其余的情况，在看完本证明后容易自行给出。记$E$上的有界、仅在有限测度集上不为0的函数为$\mathcal{F} $。

    一方面，$\forall h,k\in \mathcal{F},0\leq h\leq f,0\leq k\leq g$，有$h+k\in \mathcal{F}$，且$0\leq h+k\leq f+g$，因此
    \[\int_{E}h\,dm+\int_{E}k\,dm=\int_{E}(h+k)\,dm\leq \int_{E}(f+g)\,dm\]
    对$h,k$取上确界，即得
    \[\int_{E}f\,dm+\int_{E}g\,dm\leq \int_{E}(f+g)\,dm\]

    另一方面，$\forall l\in \mathcal{F},0\leq l\leq f+g$，定义$h=\min\{f,l\},k=l-h$，则$0\leq h\leq f,0\leq k\leq g$，其中后一不等式可通过反证法验证：
    \[g<k=l-h\implies l>g+h\geq g+f,\text{矛盾}\]
    因此
    \[\int_{E}l\,dm=\int_{E}h\,dm+\int_{E}k\,dm\leq \int_{E}f\,dm+\int_{E}g\,dm\]
    对$l$取上确界，即得
    \[\int_{E}(f+g)\,dm\leq \int_{E}f\,dm+\int_{E}g\,dm\]
\end{proof}
类似于有界可测函数，非负可测函数的勒贝格积分也具有单调性和区域可加性。

下面的命题表明，如果不涉及具体的极限过程，\textbf{总有$0\cdot\infty=0$}：

\begin{proposition}[$0\cdot\infty=0$]
    设$f:E\to[0,+\infty)$是非负可测函数，则\\
    \ding{172}$\int_{E}f\,dm=0\iff f=0,\text{a.e.}$；\\
    \ding{173}若$E_0$是$E$的零测子集，则$\int_{E}f\,dm=\int_{E\setminus E_0}f\,dm$
    %事实上，对广义实值的非负可测函数$f:E\to[0,+\infty]$（例如定义$f(x)=1/x,x>0$，则自然可以定义$f(0)=+\infty$），仍然可以定义勒贝格积分。此时：\\
    %\ding{174}如果$f$可积，则$f$几乎处处有限。

    %注意几乎处处有限不代表几乎处处有界。
\end{proposition}

特别地，单点集是零测集，因此考虑区间上的勒贝格积分时，区间的开闭总是没有影响的，因此，以下记号是良定义的：
\[\int_{a}^{b}f\,dm=\int_{[a,b]}f\,dm=\int_{(a,b)}f\,dm=\int_{[a,b)}f\,dm=\int_{(a,b]}f\,dm\]
我们当然也承认$\int_{a}^{b}f(x)\,dx$的记号。

下面引入分析学中的两个重要工具。
\begin{theorem}[Fatou引理]
    设$\{f_n:E\to[0,+\infty)\}$是非负可测函数列，$f_n\to f,\text{a.e.}$，则
    \[\int_{E}f\,dm\leq \liminf_{n\to\infty}\int_{E}f_n\,dm\]
    此定理也可表述为
    \[\int_{E}\liminf_{n\to\infty}f_n\,dm\leq \liminf_{n\to\infty}\int_{E}f_n\,dm\]
\end{theorem}

之所以称之为引理，是因为在后续的一些列定理均需要借助它来完成，但它的重要性完全不亚于其他的定理（当然，既然我们略去了多数证明，本书中没有使用Fatou引理的例子）。

\begin{theorem}[Levi单调收敛定理]
    设$\{f_n:E\to[0,+\infty)\}$是递增的非负可测函数列，且$f_n\to f,\text{a.e.}$，则
    \[\int_{E}f\,dm=\lim_{n\to\infty}\int_{E}f_n\,dm\]
\end{theorem}
后文中将它简称为“单调收敛定理”。它表明对于非负递增的的可测函数列，总能\textbf{交换极限与积分}的顺序。特别地，对于非负可测的函数项级数$\sum_{n=1}^{\infty}f_n$，其部分和列$\{s_n=\sum_{i=1}^{n}f_i\}$递增，考虑其部分和$s_n$，如果$s_n\to s,\text{a.e.}$，则
\[\int_{E}s\,dm=\lim_{n\to\infty}\int_{E}s_n\,dm=\lim_{n\to\infty}\int_{E}\sum_{i=1}^{n}f_i\,dm=\sum_{n=1}^{\infty}\int_{E}f_n\,dm\]
因此，非负可测函数项级数可以\textbf{逐项积分}。

本小节的最后，我们通过非负可测函数来定义任意可测函数的勒贝格积分。

\begin{definition}[函数的正部与负部]
    函数$f:E\to\mathbb{R}$的\textbf{正部}与\textbf{负部}分别定义为$$f^+(x)=\max\{f(x),0\},f^-(x)=-\min\{f(x),0\}$$
    易见$f=f^+-f^-,|f|=f^++f^-$。
\end{definition}
后文中将定义函数的翻转为$f^-(x)=f(-x)$，请读者注意它们的区别，正部与负部这一概念仅在本节中使用。

\begin{definition}[可测函数的勒贝格积分]
    设$f:E\to\mathbb{R}$是可测函数，如果$f^+,f^-$均可积，则称$f$是可积的，并定义$f$在$E$上的勒贝格积分为
    \[\int_{E}f\,dm=\int_{E}f^+\,dm-\int_{E}f^-\,dm\]
    $f$的可积性等价于$|f|$的可积性，因为
    \[\int_{E}|f|\,dm=\int_{E}f^+\,dm+\int_{E}f^-\,dm<\infty\]
\end{definition}

由非负可测函数积分的非负性，$f$的正部、负部的积分值均非负，因此
\[\int_{E}f\,dm\leq \int_{E}|f|\,dm\]
此即勒贝格积分的绝对值不等式。

对于黎曼积分，绝对可积性强于可积性，并称可积却不绝对可积的函数是\textbf{条件可积}的，例如，\ref{sub:int_sin}中介绍的采样函数就是条件可积的。但是，从前一定义可以看出，对于勒贝格积分，\textbf{可积性与绝对可积性是等价的}。

类似于非负可测函数，可测函数的勒贝格积分也具有线性性、单调性和区域可加性。

\begin{theorem}[比较判别法]
    设$f,g:E\to\mathbb{R}$是可测函数，$g\geq|f|$，如果$g$可积，则$f$可积，并且
    \[\left|\int_{E}f\,dm\right|\leq \int_{E}|f|\,dm\leq \int_{E}g\,dm\]
\end{theorem}

下面的定理是本书中\textbf{用于极限与积分换序的主要定理}：
\begin{theorem}[Lebesgue控制收敛定理]
    设$\{f_n:E\to\mathbb{R}\}$是可测函数列，$f_n\to f,\text{a.e.}$，如果存在一个可积函数$g:E\to[0,+\infty)$，使得对于所有$n$和$x\in E$，有$|f_n(x)|\leq g(x)$，则$f$可积，并且
    \[\int_{E}f\,dm=\lim_{n\to\infty}\int_{E}f_n\,dm\]
\end{theorem}

后文中将它简称为“控制收敛定理”。\textbf{此定理容易推广到实变量的复值函数列}，只要分别考虑实部和虚部即可。

作为使用控制收敛定理的例子，也作为勒贝格积分的优秀性质，我们给出以下定理：
\begin{theorem}[勒贝格积分的可数可加性]
    设$f:E\to\mathbb{R}$是可测函数，$\{E_k\}_{k=1}^{\infty}$是可测划分，则
    \[\int_{E}f\,dm=\sum_{k=1}^{\infty}\int_{E_k}f\,dm\]
\end{theorem}
\begin{proof}
    记$\chi_n$为区间$[-n,n]$的特征函数，则$f\cdot\chi_n\to f$，且$|f\cdot\chi_n|\leq |f|$，因此
    \[\int_{E}f\,dm=\lim_{n\to\infty}\int_{E}f\cdot\chi_n\,dm=\lim_{n\to\infty}\sum_{k=1}^{\infty}\int_{E_k}f\cdot\chi_n\,dm=\sum_{k=1}^{\infty}\int_{E_k}f\,dm\]
\end{proof}
\begin{corollary}
    设$f:E\to\mathbb{R}$可积，则\\
    \ding{172}若$\{A_n\}_{n=1}^{\infty}$是上升的$E$的可测集子族，即$A_n\subset A_{n+1}\subset E,\forall n$，则
    \[\int_{\bigcup_{n=1}^{\infty}A_n}f\,dm=\lim_{n\to\infty}\int_{A_n}f\,dm\]
    \ding{173}若$\{B_n\}_{n=1}^{\infty}$是下降的$E$的可测集子族，即$B_n\supset B_{n+1}\supset E,\forall n$，且$\bigcap_{n=1}^{\infty}B_n=\emptyset$，则
    \[\lim_{n\to\infty}\int_{B_n}f\,dm=0\]
\end{corollary}
\begin{proof}
通过一个常见的手法，我们将\textbf{具有单调性的集合列转化为不交的集合列}。令$E_1=A_1,E_n=A_n\setminus A_{n-1},n\geq 2$，则$\{E_n\}_{n=1}^{\infty}$是不交的，并且$\bigcup_{n=1}^{\infty}A_n=\bigcup_{n=1}^{\infty}E_n$，由积分的可数可加性，
\[\int_{\bigcup_{n=1}^{\infty}A_n}f\,dm=\sum_{n=1}^{\infty}\int_{E_n}f\,dm=\lim_{n\to\infty}\sum_{i=1}^{n}\int_{E_i}f\,dm=\lim_{n\to\infty}\int_{\bigcup_{i=1}^{n}A_i}f\,dm=\lim_{n\to\infty}\int_{A_n}f\,dm\]
因此\ding{172}得证，考虑$B_n=E\setminus A_n$并用德摩根律：
\[\bigcup_{n=1}^{\infty}B_n=\bigcup_{n=1}^{\infty}(E\setminus A_n)=E\setminus \bigcap_{n=1}^{\infty}A_n=E\]
即可证明\ding{173}。
\end{proof}

特别地，考虑上升的集合列$\{[-n,n]\}_{n=1}^{\infty}$及$\mathbb{R}$上的可积函数$f$，有
\[\int_{\mathbb{R}}f\,dm=\lim_{n\to\infty}\int_{[-n,n]}f\,dm\]
这个性质也是黎曼积分不具备的。黎曼积分的无穷限反常积分需要考虑积分上下限是否有关联，例如主值意义下的积分有时不等于普通意义下的积分，而勒贝格积分则不存在这样的情况。

\begin{example}
    考虑函数$f(x)=x$，它在$\mathbb{R}$上不可积，但在区间$[-n,n]$上可积，并且\[\int_{[-n,n]}f\,dm=0\]
    我们定义$f$在主值意义下的积分为
    \[\text{p.v.}\int_{\mathbb{R}}f\,dm=\lim_{n\to\infty}\int_{[-n,n]}f\,dm=0\]
    在主值意义下，$f$是可积的。但在勒贝格积分的意义下，$f$不可积，因为
    \[\int_{\mathbb{R}}|f|\,dm=\int_{\mathbb{R}}|x|\,dm=+\infty\]
\end{example}

\textbf{黎曼积分的计算技巧在多数情况下同样适用于勒贝格积分}：
\begin{theorem}[勒贝格积分的变量代换]
    设$f:[a,b]\to\mathbb{R}$可积，递增函数$g:[c,d]\to \mathbb{R}$可导，并且$g(c)=a,g(d)=b$，则
    \[\int_{a}^{b}f\,dm=\int_{c}^{d}f\circ g\cdot g'\,dm\]
\end{theorem}

\begin{theorem}[勒贝格积分的分部积分法]
    设$f,g:[a,b]\to\mathbb{R}$可导，则
    \[\int_{a}^{b}f g'\,dm=\evalat{fg}{a}{b}-\int_{a}^{b}f' g\,dm\]
\end{theorem}

\subsection{常用的线性赋范空间}\label{sub:Lp_space}

本小节主要介绍两类常用的线性赋范空间：$L^p$空间和$l^p$空间$(1\leq p\leq \infty)$，前者是处理连续信号的，而后者是处理离散信号的。

\begin{definition}[$L^p$空间]
    设$E$是可测集，$L^p(E)(0<p<\infty)$空间表示在$E$上$p$次勒贝格可积的函数组成的函数空间，即
    \[L^p(E):=\left\{f:E\rightarrow \mathbb{C} \left| \int_{E}|f(t)|^p\,dm<\infty\right.\right\}\]
\end{definition}

\begin{definition}[$L^{\infty}$空间]
    设$E$是可测集，$L^{\infty}(E)$空间表示在$E$上\textbf{本性有界}(essentially bounded)的函数组成的函数空间，即
    \[L^{\infty}(E):=\left\{f:E\rightarrow \mathbb{C} \left|\exists M<\infty,|f|\leq M,a.e.\right.\right\}\]
\end{definition}

\begin{definition}[$l^p$空间]
    设$E$是可数集，$l^p(E)(0<p<\infty)$空间表示在可数集$E$上$p$次可和的函数组成的函数空间，即
    \[l^p(E):=\left\{f:E\rightarrow \mathbb{C} \left| \sum_{t\in E}|f(t)|^p<\infty\right.\right\}\]
\end{definition}

一般而言，在考虑$l^p$空间时，可数集$E$取为$\mathbb{N}$，$\mathbb{N} ^*$（数列）或$\mathbb{Z}$（离散信号）。本小节的后续内容中，默认取$E=\mathbb{N}$。

\begin{definition}[$l^{\infty}$空间]
    $l^{\infty}$表示有界数列空间，即
    \[l^{\infty}:=\left\{\{x_n\}_{n=0}^{\infty} \left|\exists M<\infty,|x_n|\leq M,\forall n\in \mathbb{N}\right.\right\}\]
\end{definition}

不难证明以下引理：
\begin{lemma}
    $\forall a,b\in\mathbb{C}$，当$p\geq 1$时，有
    \[|a+b|^p\leq 2^{p-1}\max\{|a|^p,|b|^p\}\leq 2^{p-1}(|a|^p+|b|^p)\]
    当$0<p<1$时，有
    \[|a+b|^p\leq |a|^p+|b|^p\]
\end{lemma}

由此可以推出以上几种空间都是线性空间。此外，在$0<p<1$时，$L^p(E)$和$l^p$空间是度量空间，它们的度量分别为
\[d(f,g)=\int_{E}|f-g|^p\,dm\]
\[d(\{x_n\},\{y_n\})=\sum_{n=0}^{\infty}|x_n-y_n|^p\]

不过，$0<p<1$的情况不是我们的主要目标，因为它们不构成范数空间。下面，我们给出以上几种空间在$1\leq p \leq \infty$时的范数定义，再通过Minkowski不等式证明这些范数满足三角不等式，从而证明它们是线性赋范空间。

\begin{definition}[$L^p$空间的范数]
    设$E$是可测集，如果$1\leq p<\infty$，定义
    \[\|f\|_p=\left(\int_{E}|f(t)|^p\,dm\right)^{1/p}\]
    如果$p=\infty$，定义
    \[\|f\|_{\infty}=\inf\{M<\infty||f|\leq M,\text{a.e.}\}\]
    即$f$的本性上确界。
\end{definition}

\begin{definition}[$l^p$空间的范数]
    设$1\leq p<\infty$，对于$\{x_n\}\in l^p$，定义
    \[\|\{x_n\}\|_p=\left(\sum_{n=0}^{\infty}|x_n|^p\right)^{1/p}\]
    对于$\{x_n\}\in l^{\infty}$，定义
    \[\|\{x_n\}\|_{\infty}=\sup_{n\in \mathbb{N}}|x_n|\]
\end{definition}

在$L^p$空间中，我们\textbf{不区分几乎处处相等的函数}，即如果$f=g,\text{a.e.}$则$f$与$g$在$L^p$空间中是等价的，记作$f\sim g$，因此，以上范数定义中的函数$f$可以理解为等价类中的任意一个代表元。在这个意义下，范数的同一性是成立的。容易验证，非负性和齐次性也是成立的。

通过实数的三角不等式，可以直接证明$p=1$和$p=\infty$情况下的范数满足三角不等式。对于$1<p<\infty$的情况，我们逐步建立以下不等式。

\begin{definition}[凸函数]
    函数$\phi:(a,b)\to\mathbb{R}$是\textbf{凸}(convex)的，如果\begin{align*}
        \forall x_1,x_2\in(a,b),\forall \lambda\in[0,1],\phi(\lambda x_1+(1-\lambda)x_2)\leq \lambda\phi(x_1)+(1-\lambda)\phi(x_2)
    \end{align*}
    如果等号成立当且仅当$x_1=x_2$，则称$\phi$是\textbf{严格凸}(strictly convex)的。
\end{definition}

凸函数的左、右导数必然存在，并且满足$\phi'(\alpha-)\leq \phi'(\alpha+)$。利用切线放缩的想法，可以证明Jensen不等式。

\begin{theorem}[Jensen不等式]
    设$\phi:\mathbb{R}\to\mathbb{R}$是凸函数，$f:[0,1]\to\mathbb{R}$和$\phi\circ f:[0,1]\to\mathbb{R}$可积，则
    \[\phi\left(\int_0^1 f(t)\,dt\right)\leq \int_0^1 \phi(f(t))\,dt\]
    当$\phi$是线性函数或者$f$几乎处处等于常数时，等号成立。
\end{theorem}
\begin{proof}
    令$\alpha=\int_0^1 f(t)\,dt$，任取$m\in[\phi'(\alpha-),\phi'(\alpha+)]$，则$\phi$的图像在点$(\alpha,\phi(\alpha))$处有支撑线（即恒在$\phi$的图像下方的直线）$y=m(x-\alpha)+\phi(\alpha)$，满足\begin{align*}
        \forall t\in\mathbb{R},\phi(t)\geq m(t-\alpha)+\phi(\alpha)
    \end{align*}
    因此\begin{align*}
        \phi(f(x))\geq m(f(x)-\alpha)+\phi(\alpha)
    \end{align*}
    在区间$[0,1]$上积分，得到\begin{align*}
        \int_0^1 \phi(f(t))\,dt & \geq m\int_0^1 (f(t)-\alpha)\,dt+\int_0^1 \phi(\alpha)\,dt \\
                   & =m\left(\int_0^1 f(t)\,dt-\alpha\right)+\phi(\alpha) \\
                   & =\phi\left(\int_0^1 f(t)\,dt\right)
    \end{align*}
\end{proof}

\begin{theorem}[Young不等式]
    设$p,q\in(1,\infty)$满足$\dfrac{1}{p}+\dfrac{1}{q}=1$，则对于任意$a,b>0$，有
    \[ab\leq \frac{a^p}{p}+\frac{b^q}{q}\]
    取等条件为$a^p=b^q$。
\end{theorem}
\begin{proof}
    $\mathrm{e}^x$是严格凸函数，有\begin{align*}
        \mathrm{e}^{\lambda u+(1-\lambda)v}  \leq \lambda \mathrm{e}^u+(1-\lambda)\mathrm{e}^v
    \end{align*}
    令$\lambda=1/p,u=p\ln a,v=q\ln b$即证。
\end{proof}

$1/p+1/q=1$的关系，我们还会多次见到。在实分析中，通常称$p,q$互为共轭，但由于容易与复数的共轭混淆，我们不做此定义。\textbf{在本小节的后续内容中，$p,q$总是指满足$1/p+1/q=1$的两个正实数，并且如无特殊说明，允许其中一者为1而另一者为$\infty$}。

\begin{theorem}[Hölder不等式]
    设$f\in L^p(E),g\in L^q(E)$，有$f\cdot g\in L^1(E)$，并且
    \[\int_{E}|fg|\,dm\leq \|f\|_p\cdot \|g\|_q\]
    取等条件为存在$\lambda>0$使得$|f|^p=\lambda |g|^q,\text{a.e.}$
\end{theorem}
特别地，在$p=q=2$的情况下，Hölder不等式退化为Cauchy-Schwarz不等式：\begin{align*}
    f,g\in L^2(E)\implies f\cdot g\in L^1(E),\int_{E}|fg|\,dm\leq \|f\|_2\cdot \|g\|_2
\end{align*}
\begin{proof}
    如果$p=1,q=\infty$，不等式显然成立，$p=\infty,q=1$的情况同理。下面我们证明$1<p,q<\infty$的情况。

    不妨设$\|f\|_p=\|g\|_q=1$，否则可以将$f,g$除以对应的范数。用Young不等式，\begin{align*}
        |f(x)g(x)|\leq \frac{|f(x)|^p}{p}+\frac{|g(x)|^q}{q}
    \end{align*}
    因此\begin{align*}
        \int_{E}|fg|\,dm&\leq \int_{E}\frac{|f|^p}{p}+\frac{|g|^q}{q}\,dm\\
        &=\frac{1}{p}\int_{E}|f|^p\,dm+\frac{1}{q}\int_{E}|g|^q\,dm\\
        &=\frac{1}{p}+\frac{1}{q}=1
    \end{align*}
\end{proof}

为了证明Minkowski不等式，我们还需要以下引理：
\begin{lemma}
    设$f\in L^p(E),g\in L^q(E)$。在$p<\infty$时，$f$有对偶函数
    $$f^*=\|f\|_p^{1-p}\cdot\text{sgn}(f)\cdot |f|^{p-1}$$
    满足：\begin{itemize}
        \item $f^*\in L^q(E)$，且$\|f^*\|_q=1$；
        \item $\|f\|_p=\int_{E}f\cdot f^*\,dm$。
    \end{itemize}
\end{lemma}

在$p=1$时，$f$的对偶函数就是$\text{sgn}(f)$。这个引理直接带入即可验证。

\begin{theorem}[Minkowski不等式]
    $\forall f,g\in L^p(E)$，有$f+g\in L^p(E)$，并且
    \[\|f+g\|_p\leq \|f\|_p+\|g\|_p\]
    取等条件为存在$\lambda>0$使得$|f|=\lambda |g|,\text{a.e.}$
\end{theorem}
\begin{proof}
    $p=1$或$p=\infty$的情况显然成立，下面我们证明$1<p<\infty$的情况。

    由前面的引理知$f+g\in L^p(E)$。用Hölder不等式，\begin{align*}
    \|f+g\|_p^p & =\int_{E}(f+g)(f+g)^*\,dm \\
                 & =\int_{E}f\cdot (f+g)^*\,dm+\int_{E}g\cdot (f+g)^*\,dm \\
                 & \leq \|f\|_p\cdot \|(f+g)^*\|_q+\|g\|_p\cdot \|(f+g)^*\|_q \\
                 & =\|f\|_p+\|g\|_p
    \end{align*}
\end{proof}

Minkowski不等式表明，$L^p(E)$空间的范数满足三角不等式，因此是线性赋范空间。

利用积分的连续性，只要令$\mathbb{R} $上的函数在区间$[n,n+1)(n\in\mathbb{Z} )$上取常值，即可得到上述不等式的离散形式，从而证明$l^p$空间也是线性赋范空间。

这个不等式还能够推出有限测度集上$L^p$空间的包含关系：
\begin{corollary}
    设$m(E)<\infty$，则$\forall 1\leq p_1<p_2\leq \infty$，有$L^{p_2}(E)\subset L^{p_1}(E)$，并且\begin{align*}
        \exists c\geq 0,\forall f\in L^{p_2}(E),\|f\|_{p_1}\leq c\cdot \|f\|_{p_2}
    \end{align*}
\end{corollary}
\begin{proof}
    $p_2=\infty$的情况显然成立，下面我们证明$1\leq p_1<p_2<\infty$的情况。

    令$p=p_2/p_1>1,1/p+1/q=1$，则\begin{align*}
        \|f\|_{p_1}^{p_1} & =\int_{E}|f|^{p_1}\,dm \\
                     & =\int_{E}|f|^{p_1}\cdot \chi_E\,dm \\
                     & \leq \left(\int_{E}|f|^{p_1\cdot p}\,dm\right)^{1/p}\cdot \left(\int_{E}\chi_E^q\,dm\right)^{1/q} \\
                     & =m(E)^{1/q}\cdot \|f\|_{p_2}^{p_1}
    \end{align*}
    即\begin{align*}
        \|f\|_{p_1}\leq m(E)^{\frac{1}{p_1}-\frac{1}{p_2}}\cdot \|f\|_{p_2}
    \end{align*}
\end{proof}

对于无限测度集，$L^p$空间互不包含。例如，\ref{sub:int_sin}中指出采样函数在$\mathbb{R} $上条件可积而不绝对可积，因此$\sin t/t\notin L^1(\mathbb{R} )$，但容易验证，$\sin t/t\in L^2(\mathbb{R} )$。

$L^p$空间的完备性证明较为复杂，我们仅仅给出结论：
\begin{claim}[Riesz-Fischer定理]
    设$E$可测，$1\leq p\leq \infty$，则$L^p(E)$空间是完备的。
\end{claim}
容易由$L^p(\mathbb{R} )$的完备性推知$l^p$空间的完备性。至此，我们证明了：\textbf{$L^p$空间和$l^p$空间都是完备的线性赋范空间，即Banach空间}。

另外，在$p=2$时，容易给出$L^2(E)$空间的内积（$l^2$空间同理）：
\begin{definition}[$L^2$空间的内积]
    设$E$可测，定义$L^2(E)$空间的内积为
    \[(f,g)=\int_{E}f\cdot \overline{g}\,dm\]
\end{definition}

为了区分于分布理论的符号，本书中总是用$(\cdot,\cdot)$来表示$L^2(E)$空间的内积。

显然，$L^2(E)$空间的范数可以由这个内积导出。因此，\textbf{$L^2(E)$空间构成完备的内积空间，即希尔伯特空间}。

%下面考虑$L^p(E)$空间的对偶空间。由Hölder不等式可知，任取$L^q(E)$中的函数$g$，$g$都定义了一个$L^p(E)$的泛函：\begin{align*}
%    \Phi(f)=\int_{E}f\cdot g\,dm,\forall f\in L^p(E)
%\end{align*}

%\begin{definition}
%    称泛函$\Phi\in [L^p(E)]^*$由函数$g\in L^q(E)$\textbf{表示}(represent)，如果\begin{align*}
%        \forall f\in L^p(E),\Phi(f)=\int_{E}f\cdot g\,dm
%    \end{align*}
%    记为$\Phi=\mathcal{R} (g)$。
%\end{definition}

%事实上，$[L^p(E)]^*$中的元素总能被$g\in L^q(E)$表示：

%\begin{claim}[$L^p(E)$的对偶的Riesz表示定理]
%    设$E$可测，$1\leq p\leq \infty$，则$\forall \Phi\in[L^p(E)]^*$，存在唯一的$g\in L^q(E)$，使得$\Phi=\mathcal{R} (g)$。
%\end{claim}

%可以认为，$\mathcal{R} $是一个从$L^q(E)$到$[L^p(E)]^*$的映射。根据上述定理，$\mathcal{R} $是同构映射，\textbf{线性赋范空间$L^q(E)$与$[L^p(E)]^*$同构}。

\subsection{正交函数系}\label{sub:orthognal}
\begin{definition}[正交基]
    如果有内积空间$V$的一组基$\{\mathbf{v}_1,\mathbf{v}_2,\dots ,\mathbf{v}_n\}$，满足$\langle \mathbf{v}_i,\mathbf{v}_j\rangle =0,i \neq j$，则称这组基是一组\textbf{正交基}(orthognal bases)。\\如果对$V$中任一向量$\mathbf{w}$，均有以下分解式：
    \[\mathbf{w}=\sum_{i = 1}^{n}  c_i \mathbf{v}_i\]
则这组基是\textbf{完备正交基}(complete orthognal bases)。\\
如果$\{\mathbf{v}_1,\mathbf{v}_2,\dots ,\mathbf{v}_n\}$还满足$\langle \mathbf{v}_i,\mathbf{v}_i\rangle =1,i\in \{1,2,\dots ,n\}$，称这组基是\textbf{标准完备正交基} (complete othornormal bases)，此时空间中任意向量均有分解式
\begin{align*}
    \mathbf{w}=\sum_{i=1}^{n}\langle \mathbf{w},\mathbf{v}_i\rangle \mathbf{v}_i
\end{align*}
\end{definition}
我们对向量在标准正交基下的分解式做一些解释。在$\{\mathbf{v}_1,\mathbf{v}_2,\dots ,\mathbf{v}_n\}$是完备正交基的情况下，就可以将$\mathbf{w}=\sum_{i = 1}^{n}  c_i \mathbf{v}_i$两边同时对$\mathbf{v}_j$做内积，得到
\begin{align*}
    \langle \mathbf{w},\mathbf{v}_j \rangle=\left\langle \sum_{i = 1}^{n}  c_i \mathbf{v}_i,\mathbf{v}_j \right\rangle =c_j\left\langle \mathbf{v}_j,\mathbf{v}_j \right\rangle \\
    c_j=\frac{\langle \mathbf{w},\mathbf{v}_j \rangle}{\langle \mathbf{v}_j,\mathbf{v}_j \rangle},
    \mathbf{w}=\sum_{j = 1}^{n}  \frac{\langle \mathbf{w},\mathbf{v}_j \rangle}{\langle \mathbf{v}_j,\mathbf{v}_j \rangle} \mathbf{v}_j
\end{align*}
而标准正交基是$\langle \mathbf{v}_j,\mathbf{v}_j \rangle=1$的情况。

给出了标准完备正交基后，就可以得到以下命题，它是勾股定理的高维推广。
\begin{proposition}[勾股定理/毕达哥拉斯恒等式]
    \[|\mathbf{w}|^2=\sum_{i=1}^{n}|\langle\mathbf{w,v}_i\rangle|^2|\mathbf{v}_i|^2=\sum_{i=1}^{n}|\langle\mathbf{w,v}_i\rangle|^2\]
\end{proposition}
直接将$\mathbf{w}$与自己做内积并将之前得到的分解式带入，很容易证明这个命题。
\begin{proof}\begin{align*}
    |\mathbf{w}|^2 & =\left\langle \sum_{i=1}^{n}\langle \mathbf{w},\mathbf{v}_i\rangle \mathbf{v}_i,\sum_{i=1}^{n}\langle \mathbf{w},\mathbf{v}_i\rangle \mathbf{v}_i\right\rangle                                           \\
                   & =\sum_{i=1}^{n}|\langle\mathbf{w,v}_i\rangle|^2\langle\mathbf{v}_i,\mathbf{v}_i\rangle+\sum_{1\leq i<j\leq n}\langle \mathbf{w,v}_i\rangle\langle \mathbf{w,v}_j\rangle\langle \mathbf{v}_i,\mathbf{v}_j\rangle \\
                   & =\sum_{i=1}^{n}|\langle\mathbf{w,v}_i\rangle|^2|\mathbf{v}_i|^2=\sum_{i=1}^{n}|\langle\mathbf{w,v}_i\rangle|^2
\end{align*}\end{proof}
可见做正交基分解能够极大地简化对线性赋范空间的研究。

我们知道，$L^2([0,T])$空间是完备的内积空间（见\ref{sub:Lp_space}），在这个空间上可以定义内积：
\begin{definition}[函数空间上的内积]
    在空间\[L^2([0,T]):=\left\{f:[0,T]\rightarrow \mathbb{C} \left| \int_{T}|f(t)|^2\,\text dt<\infty\right.\right\}\]
上定义内积：
\[(f,g):=\int_{T}f(t)\overline{g(t)}\,dt \]
\end{definition}

有了内积，就可以定义正交性，从而有正交函数系的概念：
\begin{definition}[正交函数系]
    在$L^2([0,T])$空间中，如果函数列$\{f_i(t)\}_{i=1}^{\infty}$满足
\[(f_i,f_j)=\int_{T}f_i(t)f_j^*(t)\,dt=0,i\neq j\]
则称$\{f_i(t)\}_{i=1}^{\infty}$为$L^2([0,T])$空间中的一组\textbf{正交函数系}。
\end{definition}

\begin{definition}[完备正交函数系]
    对于一个正交函数系$\{f_i(t)\}_{i=1}^{\infty}$，如果不存在非零的函数$g(t)\notin\{f_i(t)\}_{i=1}^{\infty}$使得$g(t)$与$\{f_i(t)\}_{i=1}^{\infty}$中的所有函数正交，就称$\{f_i(t)\}_{i=1}^{\infty}$为\textbf{完备正交函数系}
\end{definition}

完备性意味着$L^2([0,T])$空间中的任一函数$f(t)$均可分解为这个函数系的级数$\sum_{i = 1}^{\infty}  a_i f_i(t)$，我们沿用前文中给出的公式计算这些系数：
\[a_i=\frac{(g,f_i)}{(f_i,f_i)}=\frac{\int_{T}g(t)f_i^*(t)\,dt}{\int_{T}|f_i(t)|^2\,dt}\]
细心的读者可能已经发现，这里得到的公式用到了有限维线性空间中的结论，但要推广到无限维线性空间并不是显然的。我们将在\ref{sub:paseval}中给出帕塞瓦尔恒等式之后一并讨论这个问题。

在适当的伸缩变换下，\ref{sub:bessel}中介绍的贝塞尔函数就是一个完备正交函数，傅里叶级数理论也是建立在完备正交函数系的基础之上的。此外，典型的标准完备正交函数系还有勒让德多项式 (Legendre polynomials)、小波变换基函数 (wavelet basis functions)等，下面仅讨论三角函数系和指数函数系。
